\chapter{Restricted Power Series}

Recall formal multivariate power series over a set $K$ in $n$ variables, written
$K [ \! [ x_1, \dots, x_n ] \! ]$ are infinite sums
\[
  f(x) = \sum a_{(i_1, \cdots, i_n)} x_1^{i_1} \cdots x_n^{i_n},
\]
where the sum is taken over all $n$-tuples $(i_1, \cdots i_n)$ with $i_j \in \mathbb{N}$ and
$a_{(i_1, \cdots, i_n)} \in K$.

When $K$ is a ring, formal multivariate power series form a ring, and we shall refer to them as the
ring of power series over $K$.  We again denote such rings using $K [ \! [ x_1, \dots, x_n ] \! ]$.

Note, henceforth we shall drop the word multivariate and just say power series for simplicity.

We define a subset of $K [ \! [ x_1, \dots, x_n ] \! ]$ that consists of power series whose coefficients
approach zero in the cofinite filer.

\begin{definition}
  Let $f$ be a power series over a nonarchimedean normed ring $R$ in $n$ variables, write
  $f(x) = \sum a_{(i_1, \cdots, i_n)} x_1^{i_1} \cdots x_n^{i_n}$, and let
  $c = (c_1, \cdots, c_n)$ be an $n$-tuple of real numbers.
  We say $f$ is a restricted power series of parameter $c$ over $R$ in $n$ variables if
  \[
    \lim_{n \to \infty} \left\lvert a_{(i_1, \cdots, i_n)} \right\rvert \prod_{j=1}^{n} c_j^{i_j} = 0
  \]
  in the cofinite topology.

  We denote the ring of restricted power series of parameter $c$ over $R$ in $n$ variables by
  $R_c \langle \! x_1, \dots, x_n \! \rangle$.
\end{definition}

The key specification of this will be when $c = 1$ and $R$ is a nonarchimedean normed field. We will
call such rings the Tate algebra and denote it as $T_n$.

Note there is an alternative definition \cite{...} %%Def from wikipedia for now... think Judiths
%% notes could be used for tis

\begin{definition}
  Let $A$ be a linearly topologised ring, seperated and complete and $\{ I_\lamda \}$ the fundamental
  system of open ideals. Then the ring of restricted power series is defined as the projective limit
  of the polynomial rings over $A / I_\lambda$:
  \[
    A \langle \! x_1, \dots, x_n \! \rangle = \varprojlim_\lambda A / I_\lambda [x_1, \dots, x_n].
  \]
\end{definition}

It is not hard to see that we can identify such rings with our initial definition since each $I_\lambda$
contains all but finitely many coefficients. We decide to formalise the definition in Lean as the former
as this will be more convenient for computations. But API saying it is equivalent to this projective
limit should be added to Mathlib eventually.

In literature, the ring of restricted power series are sometimes only defined for non-negative $c$-tuples,
\cite{Gouvea}, however we note by the following lemma that this does not matter and keeping with
Mathlib generality we allow any $c$.

\begin{lemma}
  A power series $f$ is restricted for parameter $c$ if and only if it is restricted for parameter
  $\lvert c \rvert$.
\end{lemma}

\section{Structure}

In the definition we call restricted power series a ring, so we better hope they are a ring. Most of
the ring axioms follow immediately from the definition of being restricted, the exceptions are
closure of the operations.
Addition follows from the usual squeeze theorem, but multiplication is of
more interest.
The nuance comes from coefficients of the product being antidiagonal sums of coefficients
of the two power series, and it may not be immediately clear that we will get the wanted restricted
condition.

To prove closure of multiplication we will use the two more general lemmas.

We first give the definition of the antidiagonal property on additive monoids.

\begin{definition} % HasAntidiagonal Algebra.Order.Antidiag
  An additive monoid, $M$ is said to have the antidiagonal property if there exists a function
  $ \phi : M \to {S | S \subset M \times M, \; S \text{ finite}}$ such that for $(s_1, s_2) \in \phi (m)$
  $s_1 + s_2 = m$.
\end{definition}

A typical example of a monoid having the antidiagonal property is $\mathbb{N}$, but it should be
immediate to see that $\mathbb{Z}$ does not. This property's importance can be seen in definitions
of Laurent series and Hahn series. %% Could explain what I mean by this


%% I think I may need a few definitions discussing how Lean handles filters and convergence
%% It is worth remarking that with Leans vast filter library we rarely need to resort to hard
%% analysis
\begin{lemma}
  Let $M$ be an additive monoid that has the antidiagonal property and $f : M \times M \to R$ be
  a function to a linearly ordered set, that has a filter $F$. If for every neighbourhood
  $a$ in $F$, the preimage of $a$ under $f$ is a neighbourhood hood in the cofinite filter,
  %% If I give a definition of tendsto I can avoid this
  then for every neighbourhood of $a$ in $F$, the primeage of $a$ under the image of the supremum of the
  antidiagonal of $a$ under $f$ is a neighbourhood in the cofinite filter.
\end{lemma}

\begin{lemma}
  Let $M$ be an additive monoid with the antidiagonal property; $S$ a nonarchimedean normed
  ring; $C : M \to \R$ be a function such that $C (a + b) = C (a) * C (b)$; and $f, g : M \to S$.
  Then if we have the functions $i \mapsto \lvert f(i)  \rvert * C(i)$ and
  $i \mapsto \lvert g(i) \rvert * C(i)$ tend to zero under the cofinite filter (that is for all
  neighbourhoods of 0 the preimage of the function is a neighbourhood in the cofinite filter). Then
  the function $a \mapsto \lvert \sum_{p \in \text{antidiagonal} a} f (p.1) g (p.2) \rvert * C (a)$
  tends to zero under the cofinite filter.
\end{lemma}

The closure of multiplication is then immediate taking $C$ to be $i \mapsto \prod c_i^i$ and $f, g$
the coefficient function mapping a power series to its coefficients, since the conditions about tending
to zero are just the definition of $f,g$ being restricted power series.

\section{Remarks}

As with many things there are decisions on typing in Mathlib, because we want to work with the ring
of restricted power series as their own thing, and not just power series that are restricted, we
first show the structure properties on power series with a predicate saying that they are restricted.
%% e.g. add the IsRestricted def here
After we show they are a subring of power series we then promote them to their own type. However,
note this means given something of type power series we have a tuple of information: a power series;
and a proof that it is restricted. Therefore, if we want to apply results on power series to the
type of restricted power series we can just apply them on the first term.
